% Image credits for the photographs and AI illustrations of Book 3 (University, Year 1).
% Machine-readable license records live in images/book3/CREDITS.md; keep the
% two files in sync. The Book 3 writer owns this file: add one \item per
% photograph (inside the itemize, by chapter) as photographs land.
% Portuguese (Brazilian) edition: translated from frontmatter/image-credits-book3.tex;
% every credit, name, dataset and licence kept exactly.
\chapter*{Créditos das imagens}

Os desenhos, os mapas e os gráficos deste livro são originais; os mapas e os
gráficos desenhados a partir de dados publicados citam sua fonte na legenda ou
nesta lista. As fotografias são usadas sob suas respectivas licenças livres, com
agradecimentos a seus autores; os links completos para cada fotografia e as URLs
das licenças estão listados em \texttt{images/book3/CREDITS.md}, no repositório
do livro.

\bigskip
Ao lado das fotografias e dos desenhos originais, algumas figuras são ilustrações
geradas por inteligência artificial, produzidas com a geração de imagens da OpenAI
a partir de prompts escritos pelos editores do livro e revisadas quanto à precisão
científica antes de sua inclusão. Elas mostram apenas cenas genéricas, nunca um
lugar real, uma pessoa real ou um espécime real. O prompt completo de cada imagem
gerada está registrado em \texttt{images/book3/ai/PROMPTS.md}, no repositório.

\bigskip
\noindent Fotografias e dados, por capítulo:
\begin{itemize}\small
  \item Cap.~1 (astronomia esférica): rastros de estrelas acima do Observatório Paranal, R.~Wesson/ESO, CC BY 4.0.
  \item Cap.~2 (telescópios, detectores e fotometria): seis segmentos do espelho primário do James Webb Space Telescope antes dos testes criogênicos, NASA/MSFC/David Higginbotham, domínio público.
  \item Cap.~3 (planetologia comparada): os planetas terrestres em escala, NASA; Vênus a partir do radar do Magellan, NASA/JPL; Marte a partir do orbitador Viking 1, NASA/USGS; Europa, NASA/JPL/DLR; Io, NASA/JPL/University of Arizona; todos de domínio público.
  \item Cap.~4 (mapas geológicos): William Smith, \emph{A Delineation of the Strata of England and Wales, with part of Scotland}, 1815; domínio público (Wikimedia Commons).
  \item Cap.~5 (seções geológicas): a discordância de Hutton em Siccar Point, PawełS, CC BY-SA 4.0 (Wikimedia Commons).
  \item Cap.~6 (métodos de campo e de laboratório): um granito biotítico em lâmina delgada entre polarizadores cruzados, Darla Sondrol, CC0 (Wikimedia Commons).
  \item Cap.~7 (minerais formadores de rochas): espécimes de olivina, quartzo, biotita, almandina e hornblenda, Robert M.~Lavinsky, CC BY-SA 3.0; labradorita, Twyla Baker, CC0; romboedros de clivagem da calcita, Giovanna Canu e Eva Santini, CC BY 4.0 (todos Wikimedia Commons).
  \item Cap.~8 (mineralogia óptica): um gabro em lâmina delgada com um polarizador e entre polarizadores cruzados, Chmee2, CC BY-SA 3.0; Auguste Michel-Lévy quando estudante, domínio público (Wikimedia Commons). As cores de interferência são calculadas a partir das funções de correspondência de cores CIE 1931 e do iluminante padrão CIE D65, \copyright{} CIE, CC BY-SA 4.0.
  \item Cap.~9 (a Terra em camadas): Andrija Mohorovičić, domínio público; Inge Lehmann em 1932, fotografia de Even Neuhaus, domínio público (Wikimedia Commons). PREM: A.~M.~Dziewonski e D.~L.~Anderson (1981), como tabulado no ObsPy.
  \item Cap.~10 (o balanço térmico da Terra): fluxo de calor global de J.~H.~Davies e D.~R.~Davies (2010), \emph{Solid Earth}, CC BY 3.0; adiabática do manto de T.~Katsura (2022), \emph{JGR Solid Earth}, CC BY.
  \item Cap.~11 (isostasia e anomalias gravimétricas): o aparelho de pêndulo de Vening Meinesz, Broodt, CC BY-SA 4.0 (Wikimedia Commons). Altitudes do ETOPO1 (NOAA National Centers for Environmental Information), amostradas a um grau.
  \item Cap.~12 (cinemática das placas sobre uma esfera): uma estação permanente de posicionamento por satélite, Alpsdake, CC BY-SA 4.0 (Wikimedia Commons). Limites de placas do PB2002 (P.~Bird, 2003; GeoJSON por H.~Ahlenius, ODC-By 1.0); linhas de costa do Natural Earth (domínio público).
  \item Cap.~13 (riftes, cristas e margens passivas): fendas abertas em Þingvellir, Islândia, Olga Ernst, CC BY-SA 4.0; lavas em almofada na Dorsal de Galápagos, NOAA Okeanos Explorer Program, CC BY 2.0 (ambas Wikimedia Commons). Profundidades do ETOPO1 (NOAA National Centers for Environmental Information) e idades do assoalho oceânico de M.~Seton e outros (2020), EarthByte, CC BY 4.0.
  \item Cap.~14 (margens ativas): o Pico Sarychev em erupção, visto da Estação Espacial Internacional, NASA, domínio público (Wikimedia Commons). Terremotos do USGS ANSS Comprehensive Catalog (domínio público); vulcões do Smithsonian Institution Global Volcanism Program; limites de placas do PB2002 (P.~Bird, 2003; GeoJSON por H.~Ahlenius, ODC-By 1.0); costas do Natural Earth (domínio público); idades do assoalho oceânico de M.~Seton e outros (2020), EarthByte, CC BY 4.0.
  \item Cap.~15 (pontos quentes e vulcanismo intraplaca): uma fonte de lava em Pu\textquoteleft u \textquoteleft O\textquoteleft o, Kilauea, G.~E.~Ulrich, U.S. Geological Survey, domínio público; os Ghats Ocidentais em Matheran, Nicholas (Nichalp), CC BY-SA 2.5 (ambas Wikimedia Commons). Posições dos pontos quentes, azimutes das trilhas e taxas da Tabela~1 do suplemento eletrônico de W.~J.~Morgan e J.~P.~Morgan (2007); profundidades do ETOPO1 (NOAA National Centers for Environmental Information); limites de placas do PB2002 (P.~Bird, 2003; GeoJSON por H.~Ahlenius, ODC-By 1.0); costas do Natural Earth (domínio público).
  \item Cap.~16 (tensão, deformação e reologia): uma célula triaxial de alta pressão para medir a resistência de rochas, Darryl Peroni, CSIRO, CC BY 3.0 (Wikimedia Commons). Parâmetros de fluência por deslocamento do quartzito úmido (Gleason e Tullis, 1995) e da olivina seca (Hirth e Kohlstedt, 2003), como compilados por M.~A.~Lopez-Sanchez em \emph{strength\_envelopes} (MPL-2.0).
  \item Cap.~17 (falhas, dobras e fábricas): um espelho de falha em Corona Heights, São Francisco, Dextersinister, CC BY-SA 4.0; brecha de falha do Keystone Thrust, Nevada, Jstuby, domínio público; dobras em chevron em Millook, Cornualha, Smalljim, CC BY 3.0; uma antiforma com clivagem plano-axial perto de Potes, Cantábria, Mikenorton, CC BY-SA 3.0 (todas Wikimedia Commons).
  \item Cap.~18 (a fonte do terremoto): terrenos e cercas separados pela ruptura de 1906 no Condado de Marin, \emph{Popular Science Monthly} vol.~69 (1906), domínio público; deslocamento do Wallace Creek ao longo da falha de San Andreas, R.~E.~Wallace, U.S. Geological Survey, domínio público (ambas Wikimedia Commons). Terremotos do USGS ANSS Comprehensive Catalog (domínio público); o mecanismo focal do terremoto de Tohoku de 2011 a partir do tensor de momento W-fase do USGS (domínio público).
  \item Cap.~19 (intemperismo e argilas): granito podre e areia de granito em Targassonne, Pireneus, Pinpin, CC~BY~2.5; bauxita sobre arenito caulinítico em Weipa, Austrália, Werner Schellmann, CC~BY-SA~2.5; cristais de caulinita ao microscópio eletrônico de varredura, U.S. Geological Survey, domínio público (todas Wikimedia Commons). Raios iônicos de R.~D.~Shannon (1976).
  \item Cap.~20 (erosão e transporte): a pluma de sedimentos do Mississippi captada pelo instrumento MODIS do Terra, 4~de março de 2018, imagem do NASA Earth Observatory de Joshua Stevens, domínio público (Wikimedia Commons).
  \item Cap.~21 (sedimentos clásticos e seus ambientes): estratificação cruzada no Arenito Aztec, Nevada, Chip Carroon, Bureau of Land Management, CC~BY-SA~2.0; o delta do rio Lena a partir do Landsat~7, U.S. Geological Survey, domínio público; um turbidito com as divisões da sequência de Bouma, Becke-Oese, Jo Weber, CC~BY-SA~3.0 (todas Wikimedia Commons).
  \item Cap.~22 (sedimentos carbonáticos, silicosos e evaporíticos): Bora Bora a partir da Estação Espacial Internacional, NASA Johnson Space Center, domínio público; uma célula de \emph{Emiliania huxleyi}, Alison R.~Taylor, University of North Carolina Wilmington Microscopy Facility, CC~BY~2.5; \emph{Orbulina universa}, H.~Takagi, K.~Kimoto, T.~Fujiki, H.~Saito, C.~Schmidt, M.~Kucera e K.~Moriya, CC~BY~4.0; diatomáceas do Estreito de McMurdo, Gordon T.~Taylor, Stony Brook University, NOAA Corps Collection, domínio público (todas Wikimedia Commons).
  \item Cap.~23 (diagênese e bacias sedimentares): sobrecrescimentos de quartzo no Arenito Potsdam, M.~C.~Rygel, CC~BY-SA~4.0 (Wikimedia Commons). Empilhamento do nível do mar de R.~M.~Spratt e L.~E.~Lisiecki (2016), NOAA NCEI World Data Service for Paleoclimatology; parâmetros de compactação de J.~G.~Sclater e P.~A.~F.~Christie (1980), como tabulados no pyBacktrack (EarthByte).
  \item Cap.~24 (a escala estratigráfica): o golden spike do GSSP do Selandiano em Zumaia, Josugoni, CC~BY-SA~4.0 (Wikimedia Commons); o limite Cretáceo--Paleógeno no Trinidad Lake State Park, Nationalparks, CC~BY-SA~2.5 (Wikimedia Commons). Idades-base e cores da Carta Cronoestratigráfica Internacional, Comissão Internacional de Estratigrafia, CC~BY~4.0.
  \item Cap.~25 (métodos de datação radiométrica): um grão de zircão em catodoluminescência, Emmanuel Roquette, CC~BY-SA~4.0; um espectrômetro de massa de ionização térmica, Radiogenic, CC0 (ambas Wikimedia Commons). A curva de calibração IntCal20 de P.~J.~Reimer e outros (2020), \emph{Radiocarbon} 62, CC~BY~4.0, como distribuída com o pacote rintcal (M.~Blaauw).
  \item Cap.~26 (estrutura da atmosfera e do oceano): recolhimento de um CTD e de uma roseta no Estreito de Bering, NASA Goddard Space Flight Center, domínio público (Wikimedia Commons). Temperatura e salinidade do World Ocean Atlas 2023 (NOAA National Centers for Environmental Information).
  \item Cap.~27 (radiação e transporte de calor): a Terra iluminada pelo sol vista do Deep Space Climate Observatory, NASA Goddard Space Flight Center, domínio público (Wikimedia Commons). Radiação, fluxos superficiais e precipitação do NCEP-NCAR Reanalysis~1 (E.~Kalnay e outros, 1996), dados fornecidos pela NOAA PSL, Boulder, Colorado, EUA, a partir de seu sítio em https://psl.noaa.gov.
  \item Cap.~28 (a circulação geral da atmosfera): um ciclone extratropical no Golfo do Alasca, NASA MODIS/LANCE, domínio público; a zona de convergência intertropical a partir de dados de nuvens do GOES-11, NASA Earth Observatory, domínio público (ambas Wikimedia Commons). Ventos e pressão ao nível do mar do NCEP-NCAR Reanalysis~1 (E.~Kalnay e outros, 1996), dados fornecidos pela NOAA PSL, Boulder, Colorado, EUA, a partir de seu sítio em https://psl.noaa.gov.
  \item Cap.~29 (a circulação do oceano): a Corrente do Golfo vista pela temperatura da superfície do mar, NOAA NESDIS, domínio público; um flutuador de perfilagem profunda sendo lançado, fotógrafo não identificado da NOAA, domínio público (ambas Wikimedia Commons).
\end{itemize}
